% ============================================================
% ch00 前言：如何使用这份解读
% ============================================================
\section{前言：如何使用这份解读}\label{sec:ch00}

\subsection{这份文档是什么}

这是一份\textbf{论文深度解读}，解读对象是本项目产出的研究论文
\textit{Sharpness Proxies are Predictors, Not Local Actuators: The Edge of
Stability Supports Monitoring, Not Local Curvature Control}
（\texttt{paper/main.tex}，以下简称``论文''）。

它不是论文的翻译，也不是入门教材。我们对你做两层假设：其一，你已经读过
\texttt{tutorial/} 里的《稳定边界（EoS）入门教程》——那里的 22 页教程
从梯度下降和损失景观讲起，一路铺垫了海森矩阵、特征值与尖锐度、
$2/\eta$ 稳定性阈值、EoS 两阶段动力学、渐进尖锐化、周期-2 振荡、
自稳定化、轨迹低维性与代理量。这些 EoS 概念在本解读中\textbf{直接使用、
不再重新定义}；如果你对某个基础概念感到生疏，请回教程对应章节复习。
其二，我们假设你\textbf{只具备深度学习背景}。
论文行文借用了大量深度学习日常训练中\emph{用不到}的行话：控制理论
（执行器、受控对象、分离原理）、数值线性代数（瑞利商、幂迭代、ARPACK）、
数理统计（置信区间判读、统计功效、配对差、去趋势相关）与动力系统（不动点、奇异摄动）。
这些术语全部收在本章第~\ref{sec:ch00-glossary}~节的``背景速查''里，
按领域分组、从深度学习已知概念出发类比解释；正文中每个此类术语
\textbf{首次出现处}也有一句括号速释，读到时不必前后翻找。

在此基础上，本解读要做三件教程没有做的事：

\begin{enumerate}
  \item \textbf{讲清楚论文的增量}：论文在 EoS 文献版图上的确切位置、
        它提出的问题为什么值得问、它的每个实验在设计上想排除什么
        替代解释；
  \item \textbf{解释数字背后的含义}：论文里成对的置信区间、配对差、
        效应量上界、去趋势 Pearson 各自说明什么、不说明什么（这些统计
        判读词的集中解释见第~\ref{sec:ch00-glossary}~节``背景速查''的
        统计组）；
  \item \textbf{呈现研究的过程性}：论文的最终叙事是干净的，但它背后
        是 \texttt{FINDINGS\_zh.md} 里 34 条发现的曲折演化——包括两次
        被后续实验推翻的中间结论。理解这个过程与理解结论同样重要。
\end{enumerate}

\subsection{论文一句话}

如果你只有三十秒：教程告诉你 EoS 现象\textbf{是什么}、为什么训练会
自发地运行在稳定边界上；这篇论文问的是下一个自然问题——
\textbf{我们能不能利用这个现象去操纵训练？}答案是干净利落的一分为二：

\begin{keybox}{论文的核心二分}
\begin{itemize}
  \item \textbf{预测端（成立）}：尖锐度可以廉价地\textbf{测量}。单次
        海森-向量积的梯度瑞利商 $\lhatg =
        \mathbf{g}^{\!\top}\!H\mathbf{g}/\mathbf{g}^{\!\top}\!\mathbf{g}$
        在 EoS 体制下去趋势（detrended，先剔除时间序列的共同趋势再做
        相关，见第~\ref{sec:ch03-stats}~节统计规范）相关 $+0.981\pm0.030$，近乎精确地追踪
        $\lmax$——一个可以部署的监测器。
  \item \textbf{控制端（结构性失败）}：尖锐度不能廉价地\textbf{操纵}。
        任何秩-1/代理驱动的单方向干预都无法把 $\lmax$ 压到 EoS 均衡
        之下——不完美阻尼把 $\lmax$ 原地钉住在 $2/\eta$，精确阻尼反而
        让 $\lmax$ 冲过阈值向上逃逸。唯有 SAM 式全局动作能够下行。
\end{itemize}
\end{keybox}

论文标题里的``预测器（predictor）''与``执行器（actuator）''借自控制
理论——研究``如何用观测与干预操纵系统''的工程学科：\textbf{执行器}
是把控制信号变成真实动作的装置（电机、阀门），\textbf{观测器}是从带噪
输出估计系统内部状态的模块。这句话的分量在于：一个系统可以拥有廉价而
准确的观测器（读得准），却不拥有成本相当的控制器 / 执行器（改得动）。
换成深度学习的语言：能便宜地\emph{读} $\lmax$，不等于能便宜地\emph{改写}
它——前者像 \texttt{no\_grad} 下读统计量的代码，后者像
\texttt{optimizer.step()} 里改写梯度的钩子。论文用一系列控制实验
（``控制实验''指控制端的干预实验，``控制''取``施加控制''义，不是
``对照实验''义）证明
EoS 正是这样的系统——这个类比贯穿全文，第九章会讲它形式化到什么程度、
在哪里只是类比（完整的控制论词汇对照表见本章``背景速查''A 组）。

\subsection{阅读路线图}

本解读的章节与论文结构大体平行，并按解读需要微调。若你对控制理论、
数值计算、统计或动力系统的行话不熟，建议先浏览本章末尾的``背景速查''
（第~\ref{sec:ch00-glossary}~节）再进入正文——不过正文中每个这类术语
首次出现处都有一句括号速释，直接顺序阅读也不会卡住：

\begin{itemize}
  \item \textbf{第一章（总览）}：论文的全景图——问题、方法、四个贡献
        的逻辑关系、一张读懂全文的``地图''。赶时间的话读完本章即可
        把握主线。
  \item \textbf{第二章（相关工作）}：论文在 EoS 文献版图上的位置，
        五条相关研究线与论文的精确分界。
  \item \textbf{第三章（实验设置）}：四类实验平台、随体制变化的
        稳定性阈值、曲率测量方法学、统计规范。这一章是后面所有
        实验数字的``度量衡''。
  \item \textbf{第四章（预测端）}：$\lmax$ 代理阶梯，从零计算量的
        范数代理到单 HVP 瑞利商；为什么监测探针应该用原始梯度而不是
        动量平滑后的更新方向。
  \item \textbf{第五至八章（控制端）}：论文的主体。第五章讲三类
        控制器（权重范数帽、$\gamma$ 帽、秩-1 plant）如何全部失败；
        第六章讲 X-SAM（第三方提出的``SAM $+$ 秩-1 顶部方向修正''
        组合优化器，Duan 等 2026）消融——逃逸贡献属于 SAM 的全局
        分量而非其秩-1 修正；第七章讲机制排除实验（谱迁移反驳、
        正交阻尼与滞后拆分、clamp-free 桥），它们与第八章的规模
        复现合称论文的``四个压力测试''；第八章讲 ResNet-18/CIFAR-100 的
        规模复现与``有动作仍被钉住''。
  \item \textbf{第九章（理论）}：自稳定帽机制的形式化——Damian 二维
        递推、阻尼后的均衡分析、命题 1 及其证明思路、五个缺口。
        全解读最重的一章，公式最多，但每个公式都配直觉解释。
  \item \textbf{第十章（讨论与局限）}：实践建议、七条局限逐条解读，
        包括``最尖的臂泛化最好''这个反直觉亮点。
  \item \textbf{第十一章（结语）}：这条研究路线对实践者、优化器
        研究者与理论工作者的分别启示。
\end{itemize}

\subsection{术语与记号约定}

本解读使用统一术语表，分四组：先是贯穿全文的\textbf{核心科学概念}
（英文原名 + 一句话解释，方便随时回查），接着是\textbf{教程未详述的
深度学习背景概念}（SAM、BatchNorm 的 $\gamma$ 等——教程以 EoS 为主线、
未及展开，本解读在正文首次需要处给出完整详解，此处各留一句话备查），
然后是论文自创的操作术语，最后是\textbf{深度学习之外的领域行话}
（第~\ref{sec:ch00-glossary}~节，单独成节）。第一、三组假定你已从
教程带来基本印象；第二组由本解读正文自负详解（详解位置在下面各条
标注）；第四组假定你从未学过对应课程。

\textbf{核心科学概念：}

\begin{itemize}
  \item \textbf{尖锐度}（sharpness）：损失景观在当前参数处的弯曲程度，
        量化为海森矩阵最大特征值 $\lmax$——值越大，局部极小值越“尖”；
  \item \textbf{渐进尖锐化}（progressive sharpening）：训练过程中
        $\lmax$ 自发增长的趋势——即使没有任何外部干预，训练本身也会把
        景观越走越“尖”（Li 等人 2022）；
  \item \textbf{周期-2 振荡}（period-2 oscillation）：$\lmax$ 越过
        $2/\eta$ 后，迭代点沿顶部特征向量方向“每两步一来回”的跳动；
  \item \textbf{自稳定化}（self-stabilization）：上述振荡经三阶泰勒项
        产生恢复力、把 $\lmax$ 拉回 $2/\eta$ 的自我调节机制
        （Damian 等人 2023）——教程定性故事的主角，本文理论章的定量
        骨架；
  \item \textbf{批尖锐度}（batch sharpness）：按当前 mini-batch 的
        损失算出的海森顶部特征值——随机体制下真正饱和在 $2/\eta$ 的
        量（Andreyev 等人 2024）；
  \item \textbf{瑞利商}（Rayleigh quotient）：形如
        $\mathbf{d}^{\!\top}\!H\mathbf{d}/\mathbf{d}^{\!\top}\!\mathbf{d}$
        的量——沿方向 $\mathbf{d}$ 的“平均曲率”，永远不超过 $\lmax$，
        二者相等当且仅当 $\mathbf{d}$ 与顶部特征向量对齐；
  \item \textbf{幂迭代}（power iteration）：反复用 $H$ 乘同一个向量、
        使其逐渐转向顶部特征向量的经典算法；
  \item \textbf{海森-向量积}（Hessian-vector product, HVP）：用双重
        反传计算 $Hv=\nabla(\mathbf{g}\!\cdot\!v)$ 的技巧——不必显式
        存出海森矩阵，是所有曲率测量的基本开销单位（教程第 3 节）；
  \item \textbf{热启动}（warm-start）：拿上一时刻的结果当下一步迭代的
        初始种子，摊薄单次测量成本；
  \item \textbf{重球动量}（heavy-ball momentum）：把历史梯度做指数
        平均后并入更新的经典动量法（论文记 $\mu$）；
  \item \textbf{平坦度--泛化联系}（flatness--generalization link）：
        “平坦的最小值泛化更好”的经典经验信条（Keskar 等人 2017），
        SAM 等方法的动机之源。
\end{itemize}

\textbf{教程未详述的深度学习背景概念}（详解位置见各条标注）：

\begin{itemize}
  \item \textbf{SAM}（Sharpness-Aware Minimization，锐度感知最小化，
        Foret 等人 2021）：每步先把参数沿归一化后的梯度方向``探出''
        固定半径 $\rho$、再用探出点处的梯度做更新的优化器——为``找
        平坦的最小值''而设计；本文控制端的逃逸基准与``全局动作''的
        代表（完整算法见第~\ref{sec:ch01}~章``概念详解''框，其下压
        $\lmax$ 的机制见第~\ref{sec:ch09}~章推论）；
  \item \textbf{BatchNorm 的仿射参数 $\gamma$}：BN 层把各通道激活
        标准化后施加的\emph{逐通道}可学习缩放向量（配对参数为平移
        $\beta$）；``$\|\gamma\|$''指把网络中全部 BN 层的 $\gamma$
        首尾拼接成长向量后取 $\ell_2$ 范数——BN 网络上替代权重范数的
        尺度 / 尖锐度代理（机制见第~\ref{sec:ch04}~章）；
  \item \textbf{归一化族 plain / BN / LN}：按``无归一化层 / 批归一化 /
        层归一化''划分的三类架构家族。归一化层改变``输出尺度由谁控制''
        ：plain 网络里权重范数直接控制激活与曲率，BN / LN 则把这条耦合
        链剪断、把尺度信息转移到自己的仿射参数上；Tier-0 廉价代理的
        可靠性因此随族而变（见第~\ref{sec:ch01}~章与第~\ref{sec:ch04}~章）；
  \item \textbf{teacher--student 回归}：固定一个随机教师网络生成带噪
        标签、再训练学生网络拟合的合成任务——数据生成真值已知、难度
        可用噪声幅度精确调节（见第~\ref{sec:ch03}~章）。
\end{itemize}

\textbf{论文自创的操作术语：}

\begin{itemize}
  \item \textbf{下逃逸}（flat-escape）：把 $\lmax$ 压到当前体制的
        稳定性阈值\textbf{以下}——论文设定的控制目标；
  \item \textbf{上逃逸}（sharp-escape）：帽被移除后 $\lmax$ 冲过
        阈值\textbf{向上}生长——与控制目标相反的失败方向；
  \item \textbf{执行器干预}（plant action）：直接改写梯度本身的
        优化器修改（如秩-1 阻尼），区别于只读不写的监测代理；
  \item \textbf{有动作仍被钉住}（pinned despite action）：执行器
        可验证地以全力触发，$\lmax$ 却仍留在跨种子噪声包络之内；
  \item \textbf{自稳定帽}（self-stabilization cap）：$2/\eta$ 阈值
        作为由周期-2 振荡维持的上限——``帽''的意象是它只能被顶开、
        不能被压低；
  \item \textbf{秩-1 阻尼}（rank-1 damping）：把梯度沿顶部特征向量
        的分量缩小 $1/d$ 的干预；
  \item \textbf{谱迁移}（spectral migration）：被论文反驳的替代
        假说——尖锐度像守恒预算一样从被阻尼方向流向未阻尼方向；
  \item \textbf{钳位门控}（clamp）：$d\gets\max(\cdot,1)$ 的自门控，
        只在高于目标时触发的保守设计；
  \item \textbf{伪收敛}（false convergence）：幂迭代过早停止落在
        次主导瑞利商上，系统性低估 $\lmax$ 的测量病态。
\end{itemize}

记号沿用教程：$\lmax$ 为损失海森矩阵的最大特征值（海森矩阵是损失对
全部参数的二阶偏导数构成的方阵，其特征值就是各正交参数方向上的损失
曲率；因为它的元素个数是参数量的平方、根本无法显式存储，所以全文的
曲率测量都用``只算 $Hv$、不存 $H$''的隐式技巧，见背景速查的
``海森--向量积''条），$\eta$ 为学习率，$\vtope$（论文记
$v_{\mathrm{top}}$ / $u_1$）为\textbf{顶部特征向量}——最大特征值
对应的特征方向，即曲率最陡的参数方向（类比 PCA 的第一主成分方向），
$\lhatg$ 与 $\lhatd$ 分别为沿梯度方向与沿更新方向的瑞利商。
\textbf{体制}（regime）一词在论文中有双重用法：一是指训练体制
（全批 GD / 动量 / 小批 SGD），二是指秩-1 plant 的两种结局
（体制 (i) 不完美阻尼钉在 $2/\eta$、体制 (ii) 精确阻尼上逃逸），
上下文可以区分。另有两组\textbf{复用记号}需留意（各节随文定义、以
各节为准）：$\alpha$ 在第五章是梯度沿 $\hat v$ 的分量
$\langle\mathbf{g},\hat v\rangle$，在第六章是 X-SAM 的固定超参
（$\alpha{=}0.2$，无关），在第九章是渐进尖锐化系数；$\gamma$ 在
第三至五章是 BatchNorm 的仿射尺度参数（$\|\gamma\|$、$\gamma$-帽），
在第九章是阻尼强度 $1-1/d$——两种 $\gamma$ 含义无关，且方法名
``$\gamma$-帽''永远指 BN 参数。其余复用记号各随文定义：$\mu$
在训练体制里是重球动量系数（$\mu{=}0.9$），在第九章是周期轨道乘子
（一律带下标，如 $\mu_y$）；$\mathbf{d}$ 在第四章是探针方向向量，
$d$ 在第五章起是阻尼除数标量（$d\gets\max(\hat\lambda/(2/\eta),1)$）；
$n$ 有时是训练样本数（附录 A 表）、有时是种子数、有时是参数总维数
（``秩-$(n{-}1)$ 正交补''）；$k$ 有时是滞后步数、有时是 rank-$k$
的秩、有时是 ARPACK 求的前 $k$ 个特征值数；$\varepsilon$ 有时是
SAM 的扰动向量、有时是第九章证明里的小偏差标量。

\subsection{背景速查：深度学习之外的术语}\label{sec:ch00-glossary}

论文标题的组织轴借自控制理论，测量靠数值线性代数，实验结论靠数理统计，
机制解释靠动力系统——这四门课的行话在深度学习日常训练中几乎用不到，
却在这篇论文里无处不在。本节按领域分组给出面向深度学习读者的速查
解释（每个术语在正文\textbf{首次出现处}还有一句更简短的括号速释，
两边内容一致）。解释刻意不追求严格，只求\emph{从你已经知道的东西
出发}：PCA、动量、优化器、训练曲线。

\subsubsection*{A. 控制理论：论文标题的词汇表}

控制理论研究``怎么用观测和干预操纵一个系统''（恒温器、自动驾驶、
化工厂）。它把一个控制系统拆成四个角色，与深度学习训练的对应如下：

\begin{table}[!ht]
\centering\small
\caption{控制回路四角色与深度学习训练的对应。论文标题``预测器
（predictor），而非执行器（actuator）''用的就是这套词汇——标题的
``预测器''对应表中``观测器''一行（论文用 predictor 称呼监测器）。}
\label{tab:ch00-control}
\begin{tabular}{lp{0.30\linewidth}p{0.34\linewidth}}
\toprule
控制论词（原义） & 在本文中 & 深度学习类比 \\
\midrule
观测器 observer（从带噪输出估计系统内部状态的模块） & 监测器：只读不写地测 $\lmax$ & \texttt{no\_grad} 下读 loss / 统计量的代码 \\
控制器 controller（决定施加多大动作的决策逻辑） & 干预规则的设计（阻尼多大、何时触发） & 更新规则的计算逻辑 \\
执行器 actuator（把控制信号变成物理动作的装置，如电机、阀门） & 直接改写梯度、介入训练的那一步 & \texttt{optimizer.step()} 里改写梯度的钩子 \\
plant 受控对象（被操纵的系统本体，工厂行话里的``装置''） & 被干预的训练过程本身。注意：论文行文把出现在\emph{方法名}里的 plant（``秩-1 plant''、``true-$\lmax$ plant''、``plant action''）借指\emph{作用于训练过程的干预模块 / 优化器变体}——读到这类名字一律按``干预''理解 & 训练循环（plant 的原义）／改写梯度的自定义优化器（方法名里的 plant） \\
\bottomrule
\end{tabular}
\end{table}

其余高频控制论词：

\begin{itemize}
  \item \textbf{传感}（sensing / sensor）：控制闭环（闭环 = 边测边调、
        测量结果反馈进下一步决策的回路；开环 = 只按预定指令行动、不看
        结果——恒温器是最简单的闭环）里为干预提供实时
        测量的环节——本文指每步估计特征向量 / 曲率的那几行代码（类比
        Adam 每步先算的一阶/二阶矩）。传感与监测的区别：传感是
        \emph{为干预服务}的测量（是控制器的一部分），监测是只读的
        观察；
  \item \textbf{阻尼}（damping）：任何``把振荡往小压''的机制，
        物理原型是摩擦 / 减震器；本文指把梯度某个方向的分量除以
        $d>1$，相当于给该方向的来回振荡加摩擦。\textbf{反向阻尼}
        （负阻尼）则相反——不压制反而放大振荡，可致发散；
  \item \textbf{增益}（gain）：放大倍数——一个系数每步把信号乘多大；
  \item \textbf{门控 / 钳位}（gate / clamp）：``门''指只在条件满足时
        才放行干预（如 $\max(\cdot,1)$：低于阈值就什么都不做，
        类比 MoE 路由器决定是否启用专家）；``界''指干预每步都发生、
        但幅度被截断在固定区间内。``自门控''指开门条件由自己测到的量
        决定；
  \item \textbf{伺服}（servo）：持续把某个量拉回目标值的闭环装置
        （恒温器是最常见的例子）；本文``幅度伺服''指振荡幅度自动
        充当恒温器、回调尖锐度；
  \item \textbf{观测器--控制器分离原理}（observer--controller
        separation principle）：线性控制论定理——``估计系统内部状态''
        与``设计干预''可以分开设计、拼起来仍然稳定。它不平凡，因为
        直觉上两者互相纠缠（控制依赖估计，估计又被控制扰动）。工程上
        它的价值是：控制器可以先按``状态完全已知''设计，再独立补一个
        观测器。论文只借它类比``测尖锐度''与``控尖锐度''是两个可以
        分开回答的问题；至于答案会不会互相牵连（例如执行器干扰自己的
        传感器），线性定理的保证在此不适用——那正是论文要用实验回答
        的部分，论文也诚实声明对非线性、非凸的网络训练这只是类比而非
        定理。
\end{itemize}

\subsubsection*{B. 数值线性代数：测曲率的工具箱}

\begin{itemize}
  \item \textbf{谱}（spectrum）：一个矩阵\emph{全部特征值}的集合。
        ``海森谱''即每个参数方向上损失曲率的完整清单，类比 PCA 中
        全部主成分的方差列表。\textbf{重尾谱}：少数特征值极大、
        大量方向仍有可观曲率的拖尾分布；
  \item \textbf{顶部 / 次主导}（top / subdominant）：
        ``顶部特征值''= 最大特征值 $\lmax$，``顶部特征向量''=
        它对应的方向（PCA 第一主成分）；``次主导''= 第二大特征值
        （及其方向）；
  \item \textbf{简并 / 近简并 / 简单特征值}（degenerate /
        near-degenerate / simple）：特征值``简单''指它只对应一个
        特征方向、没有并列第一；``（近）简并''指（几乎）有多个方向
        共享同一特征值——此时``谁是最尖方向''（几乎）分不出伯仲，
        正是幂迭代最容易出错的谱形；
  \item \textbf{特征间隙 / 谱缺口}（eigengap / spectral gap）：
        最大与第二大特征值之差 $\lambda_1-\lambda_2$。间隙越大，
        ``最尖方向''越独一无二，幂迭代越快越稳（类比：PCA 前两个
        主成分方差差得越远，第一主成分越稳定）；
  \item \textbf{瑞利商}（Rayleigh quotient）：
        $R(\mathbf{d})=\mathbf{d}^{\!\top}\!H\mathbf{d}/
        \mathbf{d}^{\!\top}\!\mathbf{d}$。把 $\mathbf{d}$ 分解到
        海森的各特征方向上，$R$ 就是各特征值按分量占比平方加权的
        \emph{平均}——所以它必然落在最小与最大特征值之间（是
        $\lmax$ 的下界）；$\mathbf{d}$ 与顶部特征向量夹角为 $\theta$
        时，顶部方向的权重只有 $\cos^2\theta$，读数大约被压到
        $\lmax\cos^2\theta$。这两条性质（下界 + 折减公式）是第四章
        全部数字的钥匙；
  \item \textbf{二次型}（quadratic form）：形如
        $\mathbf{x}^{\!\top}\!H\mathbf{x}$ 的标量（向量夹在矩阵
        中间）；``局部二次型''指把损失面在当前参数附近近似成二次
        函数（泰勒展开截到二阶）；
  \item \textbf{幂迭代}（power iteration）：反复做
        $\mathbf{v}\gets H\mathbf{v}/\|H\mathbf{v}\|$。最大特征方向
        的分量每次被放大得最多，若干次后 $\mathbf{v}$ 就对准顶部
        特征向量——PCA 求第一主成分用的就是这一招。收敛快慢由间隙
        决定：每迭代一次，误差大约缩小 $\lambda_2/\lambda_1$ 倍；
  \item \textbf{海森--向量积 / Pearlmutter 技巧}（Hessian-vector
        product, HVP）：先反传得梯度 $\mathbf{g}$，再对标量
        $\mathbf{g}\!\cdot\!\mathbf{v}$ 做一次反传即得 $H\mathbf{v}$
        ——即 PyTorch 里对反传结果再开计算图反传一次
        （\texttt{create\_graph=True}）。因为海森有参数量平方个
        元素、存不下，这种``只算 $Hv$、不存 $H$''的技巧是全部
        曲率测量的基本开销单位；
  \item \textbf{热启动}（warm-start）：拿上一时刻的结果当下一步
        迭代的初始种子，摊薄单次测量成本；
  \item \textbf{ARPACK}：经典的迭代特征值求解库（scipy 的
        \texttt{eigsh} 即其接口），基于 Lanczos 类算法——与幂迭代
        同属``只用矩阵 $\times$ 向量''的迭代法，但 Lanczos 同时维护
        一个 $k$ 维候选子空间而非单个向量，一次分清前几名特征值的
        大小次序，因此不会像幂迭代那样把第二名错当第一；它还自带
        Ritz 残差自检（Ritz 见下条），收敛不了会继续算。只要能做
        ``矩阵 $\times$ 向量''它就能求前 $k$ 个最大特征值——正好
        适配只能做 HVP、不能显式存储的神经网络海森。参数
        $k{=}3$、tol $10^{-6}$ 即``求前 3 个最大特征值、收敛容差
        百万分之一''。本文把它当作精确基准（ground truth 尺子）——
        它虽是迭代法，但 tol 压到 $10^{-6}$ 后自身误差比要测的现象
        小几个数量级（解读第
        \ref{sec:ch03-curvature} 节的对账即以它为基准）；
  \item \textbf{Ritz 值 / Ritz 残差}（Ritz value / residual）：
        迭代法给出的\emph{近似}特征值 / 特征向量（半路上的当前最佳
        猜测）；Ritz 残差
        $\|H\hat{\mathbf{v}}-\hat\lambda\hat{\mathbf{v}}\|$（除以
        $|\hat\lambda|\,\|\hat{\mathbf{v}}\|$——即除以主项的大小，
        得到不带单位的纯比值（相对误差），如 $10^{-6}$ 即``偏差只有
        主项的百万分之一''）是把这对猜测代回定义式后剩下的偏差——检验
        ``到底像不像真特征对''的误差指标，越小越可信；
  \item \textbf{预条件化 / 预条件子}（preconditioning /
        preconditioner）：用某个矩阵 $P$ 重缩放梯度
        （$P\nabla L$）再更新，等价于在拉伸过的坐标系里做梯度下降；
        Adam 的逐坐标缩放就是一种预条件。``预条件化的特征值''指在
        这个缩放坐标系里测的曲率——所以 Adam 体制下饱和的量与
        阈值都和 SGD 不同；
  \item \textbf{谱主体}（spectral bulk）：特征值分布里的``大部队''
        ——除前 $k$ 大之外的众多小特征值及其方向；
  \item \textbf{线搜}（line search）：沿一个固定方向试几个不同步长
        （各做一次前向）、挑损失最小的那个——经典优化里调步长的
        标准做法，现代深度学习训练几乎不用，本文借它当测量工具；
  \item \textbf{Davis--Kahan 定理}：经典矩阵扰动结果——特征向量的
        估计误差由``扰动大小 $/$ 特征间隙''控制；大意是前两名特征值
        差得够远时，顶部方向对噪声稳健；
  \item \textbf{Lipschitz 海森}：假设海森矩阵自身随参数的变化速率
        有上界——``曲率地图是平滑的''，分析特征向量随训练漂移的
        速度时需要的标准假设；
  \item \textbf{Fisher 矩阵}：与海森同族的``曲率 / 敏感度''矩阵
        （损失对参数扰动的统计平均敏感度），自然梯度与 K-FAC 用的
        就是它。
\end{itemize}

\subsubsection*{C. 数理统计：读懂数字的判读规则}

\begin{itemize}
  \item \textbf{置信区间}（confidence interval, CI）：对\emph{均值}
        的不确定度。``mean$\pm$95\% CI''的意思是：同样的实验若
        重复很多次，每次这样算出的区间约有 95\% 会罩住真实的平均值。
        它回答``均值有多准''而非``单次运行波动多大''（后者是标准差）
        。注意它并非天然比标准差窄：半宽约是
        $t$ 系数（见下条）$\times s/\sqrt{n}$，种子数要多（$n\gtrsim8$）才比
        单种子标准差 $s$ 窄——本文 $n{=}5$ 时半宽约 $1.24\,s$，
        比单种子波动带还宽一点，这是小样本的诚实代价；
  \item \textbf{Student-$t$ 区间}：种子数很少（本文 $n{=}5$）时，
        连``波动本身多大''都只能估，估的过程引入额外不确定性；
        Student-$t$ 分布就是小样本修正。区间半宽
        $t_{0.975,n-1}\,s/\sqrt{n}$ 中：$s$ 是种子值的样本标准差，
        $s/\sqrt{n}$ 是均值的标准误，$t_{0.975,4}\approx2.78$ 是
        查表系数（代替大样本正态的 1.96）。\textbf{保守}指区间因此
        更宽——好处是不易把纯噪声误报成真差异；\textbf{双侧}指
        两端各留 2.5\%（同时提防偏高偏低两种错判）；
  \item \textbf{CI 跨零}：对相关系数或差值而言，0 代表``没有效应''
        ；区间罩住 0 就意味着无法排除``其实毫无关系''——此时单看
        某个种子算出的 $+0.55$ 只是运气。\textbf{CI 不重叠}是偏
        严格的差异判据：两组区间完全错开时，``两组真实均值相同''
        很难成立，差异可以当结论。\textbf{上 / 下置信界}即 CI 的
        右 / 左端点；``A 的上界仍低于 B 的下界''= A 最好情形仍比
        B 最差情形低，是比均值比较更强的分离证据；
  \item \textbf{Pearson 相关系数}（Pearson correlation
        coefficient）：两条等长序列\emph{线性}同步程度的度量
        $r=\sum_i(x_i-\bar x)(y_i-\bar y)\big/
        \bigl[\sqrt{\sum_i(x_i-\bar x)^{2}}\,
        \sqrt{\sum_i(y_i-\bar y)^{2}}\bigr]$——把两条序列各自去均值后
        视为向量、取夹角余弦（无量纲，可跨序列比较）。取值 $[-1,1]$：
        $+1$ 完全同涨同跌、$0$ 涨落无线性关系、$-1$ 完全反向。本文一切
        ``相关 $+0.98$''式的数字均指它，且一律配合下条的去趋势使用
        （完整展开见第~\ref{sec:ch03-stats}~节）；
  \item \textbf{去趋势 / 滚动均值}（detrending / rolling mean）：算
        相关前的预处理——每个点减去其邻近 $w$ 个点的均值（本文
        $w{=}5$，``滚动均值 window 5''即此：慢变化趋势的局部估计），
        只留涨落\emph{残差}再做相关。必要性：两条同向单调漂移的曲线
        即使涨落彼此无关，原始 Pearson 也接近 $+1$——共同趋势制造的
        假相关；去趋势后问的才是``涨落是否同步''。展开解释与袖珍
        反例见第~\ref{sec:ch03-stats}~节。
  \item \textbf{统计不可分}：现有数据无法把两臂区分开（CI 互相
        重叠）——它是``测不出差异''，\emph{不是}``证明了没有差异''
        ；
  \item \textbf{配对差}（paired difference）：对每个种子分别算
        两臂之差、再对这 $n$ 个差值做统计。因为两臂共享同一种子的
        难度与数据噪声，相减后这些共同波动被约掉、只剩下优化器
        本身的差异——同样的种子数能分辨更小的效应；
  \item \textbf{效应量 / 效应量上界}（effect size / upper bound）：
        效应量 = 想检测的差异本身的大小（这里是两臂晚期 $\lmax$
        之差）；其置信区间的上端点即``效应量上界''——真实差异再大
        也大不过它。``上界 $+1.85$''可翻译成``就算有效应，也小于
        SAM 晚期绝对 $\lmax$ 的 $8\%$''；
  \item \textbf{统计功效}（statistical power）：当真实效应存在时，
        实验能把它检测出来的概率。种子越少、噪声越大，功效越低，
        只能检出大效应——所以``在该功效下检测不到''\emph{不等于}
        ``效应为零''。\textbf{null 结果}（阴性结果）即``在给定功效
        下未检出差异''的结论：它有信息量（排除了较大的差异），
        但不是零效应的证明。\textbf{功效分析}：实验前先计算``要
        检出预期大小的效应至少需要几个种子''的统计设计；
  \item \textbf{预注册}（pre-registered）：跑实验之前把``怎样算
        成功''的量化判据与门槛公开写死，防止看完结果再临时改标准
        、把噪声解读成胜利；
  \item \textbf{臂}（arm）：并行实验中的一个分组 / 变体（来自临床
        试验行话）——``四臂对比''即同一设置下同时跑四种配置。
        注意本文的``控制臂''指\emph{开启干预}的实验组，对照组是
        SGD 基线臂（``控制''取``施加控制''义，不是``对照''义）；
  \item \textbf{混杂变量}（confound / confounded）：同时影响干预
        与结果、会让归因出错的第三因素；好的实验设计把它固定或拆开
        检验；
  \item \textbf{备择假设}（alternative hypothesis）：与主解释竞争
        的另一种解释；``排除备择''= 设计实验证明竞争解释不成立，
        从而反证主解释；
  \item \textbf{外部效度}（external validity）：结论能否推广到
        实验设置之外（更长训练、别的数据集）；与之相对的内部效度
        指实验本身做得干不干净（本文的确定性协议就是在保内部效度）
        ；
  \item \textbf{噪声地板 / 噪声包络}（noise floor / envelope）：
        噪声地板 = 测量系统自身抖动的下限（本文同种子重跑可差约 6
        单位 $\lmax$，小于它的差异测不到）；噪声包络 = 多种子曲线
        叠在一起围出的波动带——一条臂的曲线若始终待在别的臂的
        包络里，就说它``没有响应''；
  \item \textbf{贝叶斯最优地板}（Bayes-optimal floor）：数据生成时
        注入的标签噪声决定了任何模型都不可能突破的测试误差下限
        ——即使完全知道数据生成真相的最优预测者（贝叶斯最优）也
        只能达到这个水平；
  \item \textbf{假阳性}（false positive）：把噪声误报成真效应；
        反面是漏报（真效应被噪声吞掉）。
\end{itemize}

\subsubsection*{D. 动力系统：第九章的词汇表}

\begin{itemize}
  \item \textbf{不动点}（fixed point）：迭代一步后原地不动的状态
        （类比 loss 收敛到的驻点）。\textbf{周期-2 轨道}：不是单点
        不动，而是两个点轮流切换、每两步一循环——类比学习率偏大时
        loss 在两个值之间来回跳。``轨道''指迭代产生的整条序列
        ；记号 $F^{2}$ 表示把更新规则 $F$ \emph{复合执行两步}
        （不是平方）——对齐两步一循环的周期后，周期-2 轨道就是
        $F^{2}$ 的不动点；
  \item \textbf{均衡 / 稳态}（equilibrium / steady state）：训练
        动力学的\emph{长期驻留模式}。注意本文的均衡不是静止点，
        而正是那条周期-2 振荡轨道；\textbf{瞬态}（transient）指
        进入稳态前的短暂过渡阶段（训练初期 $\lmax$ 冲上万的尖峰
        就是瞬态，不代表长期行为）；
  \item \textbf{翻号}（sign-flip）：迭代每步改变符号（$x\to-x$）
        。注意与``相关由正转负''完全无关——那是相关系数变号的
        日常用语，本文写作``由正转负''以示区分；
  \item \textbf{线性稳定性分析}：在均衡点附近把更新规则做一阶
        （线性）泰勒展开，检查扰动会被逐步放大还是缩小——对二次
        损失上的梯度下降，它给出``$\eta\lambda>2$ 就振荡发散''
        这条线；动量把这条线抬到 $2(1{+}\mu)/\eta$；
  \item \textbf{Jacobian}：更新规则的一阶导数矩阵（小扰动每步被
        如何放大 / 旋转的局部线性近似）。其特征值的模 $>1$ 表示
        扰动被放大、$<1$ 表示被吸住；
  \item \textbf{吸引子 / 排斥子}（attractor / repeller）：附近轨道
        分别被吸住 / 被推离的状态。说某轨道``是排斥子''并不妨碍
        实验里观察到它——真实轨迹可以在排斥轨道\emph{附近}贴着
        走（被慢方向的收缩拴住）而不严格锁进去；
  \item \textbf{乘子}（multiplier）：给周期轨道一个小扰动，绕轨道
        一整圈后扰动被乘的那个系数（每个方向各有一个）。$|\mu|<1$
        = 该方向扰动每圈缩小（收缩），$>1$ = 放大；
  \item \textbf{慢--快 / 奇异摄动}（slow--fast / singular
        perturbation）：系统里同时有每步剧变的快变量与每步只挪一点
        的慢变量时的标准分析方法——近似认为快变量瞬时到达均衡，
        把耦合问题拆成两个一维问题（类比：把每个 step 的高频震荡
        与几千步的缓慢漂移分成两层分析）。\textbf{快流形}：把慢
        变量暂时视为常数时快变量走出的那条轨道；\textbf{役使}
        （slaved）：快变量被慢变量锁死、只能瞬时跟随而无独立动力学
        ——``$x$ 被役使''即 $x$ 的振幅由慢变量 $y$ 当前值决定；
  \item \textbf{不变图 / 不变集}（invariant graph / set）：迭代
        映射下进入后就不再离开的点集——这里就是
        $x\in\{\pm x^{*}\}$ 两个点来回跳的集合。``图''是动力系统
        对``由状态参数化的集合''的称呼，与神经网络计算图无关；
  \item \textbf{中心流形}（center manifold）：系统里既不快速收缩
        也不快速发散的``中性''方向构成的低维集合，轨迹常被吸在
        上面慢慢演化（可类比 EoS 轨迹被 PCA 第一主成分吸住）；
  \item \textbf{分岔}（bifurcation）：参数（如学习率）连续微小
        变化时，系统行为在某个临界值处突然换模式（收敛变振荡等）
        ；分岔论即研究这种质变的分支；
  \item \textbf{相 / 相图}（phase / phase diagram）：``相''= 参数
        空间中行为模式相同的区域（类比水的固 / 液 / 气三态）；
        相图 = 标出``哪片参数区域对应哪种动力学''的地图（类比 lr
        finder 里的``太小不收敛 / 太大发散 / 中间健康''分区）；
  \item \textbf{ansatz / 试解}：先猜解的形状（如``每步变号、幅度
        不变''），代回方程验证是否自洽——类似解差分方程先假设
        指数形式再定系数；
  \item \textbf{自洽}（self-consistent）：解的取值反过来支撑解自身
        成立——两个方程互相咬合时唯一同时成立的取值；
  \item \textbf{二维约化 vs 轨迹定理}：``二维约化''指把百万维训练
        动力学压缩成两个组合变量（沿顶部方向位移 $x$、尖锐度超出
        阈值的量 $y$）再分析的降维近似——因为 EoS 行为几乎全部
        发生在顶部一个方向上（轨迹 PCA 也支持），两个变量足以
        描述核心机制；``轨迹定理''指对真实高维轨迹（而非约化
        模型）成立的定理——本文只承诺前者；
  \item \textbf{动力系统杂项}：周期-1 = 走一步即回到原处（就是不动点
        本身）；``退化''不动点 = 振幅为零的平凡 / 临界情形（与特征值
        的``简并''重名不同义）；良定义（well-defined）= 无歧义地
        存在（如 $\lambda_1$ 无并列时尖锐度才是参数的光滑函数、梯度
        才存在）；单参数形变 = 同一系统拧一个旋钮（如阻尼强度
        $\gamma$）得到的连续变形版本；物理性 = 数学解必须满足的现实
        约束（如最大特征值非负）。
\end{itemize}

\subsubsection*{E. 物理与经典优化杂项}

\begin{itemize}
  \item \textbf{重球动量}（heavy-ball momentum）：动量法的学名
        （Polyak 1964），意象是一个有质量的球滚过损失面、惯性使它
        冲过窄谷——就是你在 PyTorch 里用的 momentum，惯性 $\mu$
        越大越容易振荡，所以``开始振荡的曲率阈值''从 $2/\eta$
        抬到 $2(1{+}\mu)/\eta$；
  \item \textbf{catapult / 弹射}：大学习率训练早期，参数沿海森最尖
        方向被``弹射''离开、直到该方向曲率降到稳定阈值之下的阶段
        ——故名弹射；它是``用谱语言谈论训练稳定性''这条思路的起点
        ；
  \item \textbf{三阶泰勒项 / 三次恢复力}：损失不对称（三阶导数
        非零）时，权重振荡得越远、遇到的有效曲率变化越大，像弹簧
        回弹一样把振荡幅度拉回平衡。注意``三次''指它源自泰勒展开
        的\emph{三阶项}（曲率对位置的变化率），对慢变量的增量恰好
        是位移的\emph{二次}式——这不是笔误；
  \item \textbf{三阶张量收缩}：损失三阶导数是一个``三维数字方块''
        （张量，每个下标对应一个参数维）；把其中两个槽位填上同一个
        向量 $\mathbf{u}$ 做内积（类似 einsum 缩并两个维度），
        得到一个普通向量。恒等式 $\nabla S=\nabla^{3}L(u_1,u_1)$
        的意思是``最大特征值的梯度 = 曲率沿顶部方向的变化率''；
  \item \textbf{各向同性}（isotropic）：各方向性质相同、无偏好方向
        ——本文的``各向同性 teacher''指构造出的合成教师在多个方向
        有对称的尖锐度（人为造出两条等深峡谷，测试锐度会不会搬家）
        ；
  \item \textbf{饱和}（saturation）：到达上限后不再上升（物理里
        电容充满的意思）——$\lmax$``饱和在 $2/\eta$''即被封顶，
        与激活函数进入平台区无关；
  \item \textbf{截断}（cutoff）：本文指 SAM 情形下被压低的稳定
        阈值上限（依赖 $\|\nabla L\|$、低于 $2/\eta$），不是截断
        误差。另注意第七章的``编辑因子 $f$ 截断到 $[-1,1)$''取
        clip 夹到区间义（对应 A 组``门控 / 钳位''里的``界''），与
        本条不同；
  \item \textbf{$O(\cdot)$ 的尺度用法}：大 $O$ 在本文常作``量级 /
        尺度''用而非算法复杂度：``$O(1/\eta)$ 时间尺度''= 该量要
        约 $1/\eta$ 步才明显变化；``每步 $O(\eta)$ 漂移''= 每步
        只挪动与学习率同量级的小幅度；$O(\varepsilon^{2})$ =
        平方级小量、可忽略；
  \item \textbf{扰动分析}（perturbation analysis）：数学上的小偏差
        级数展开方法（把``每 $N$ 步才刷新一次''视为对理想情形的
        小偏差、逐阶算出影响），与 SAM 给参数加``扰动''无关。
\end{itemize}

\subsection{与项目其他文档的关系}

\begin{itemize}
  \item \texttt{tutorial/}——入门教程（本解读的前置）；
  \item \texttt{paper/main.tex}——论文原文（解读对象）；
  \item \texttt{FINDINGS\_zh.md}——34 条发现的完整演化日志，含被
        自我推翻的中间结论；
  \item \texttt{README.md}——项目总览与快速上手；
  \item \texttt{figures/}——论文原图，本解读插图全部取自于此。
\end{itemize}

本解读中的所有实验数字均逐字核对自论文原文；论文与日志不一致处
（极少，且论文为准）在解读中以论文为准。
